\chapter{Besoins fonctionnels}

L'application devra permettre de réaliser les principales opérations liées à la gestion des fichiers.

\section{Gestion des utilisateurs}

\begin{itemize}
    \item Authentifier les utilisateurs.
    \item Permettre à chaque utilisateur d'accéder à son espace personnel.
    \item Restreindre l'accès aux fichiers selon le compte utilisateur.
    \item Permettre à l'administrateur de gérer les comptes.
\end{itemize}

\section{Gestion des fichiers}

\begin{itemize}
    \item Importer des fichiers Excel compatibles.
    \item Afficher la liste des fichiers précédemment importés.
    \item Ouvrir un fichier depuis l'espace personnel.
    \item Consulter les différentes feuilles d'un classeur.
    \item Enregistrer les modifications apportées au fichier.
    \item Télécharger le fichier Excel.
    \item Télécharger une version modifiée du fichier.
    \item Permettre à l'utilisateur de gérer ses fichiers depuis son espace personnel.
\end{itemize}

\section{Édition des fichiers}

L'application devra proposer un ensemble de fonctionnalités permettant de travailler directement sur les fichiers Excel depuis le navigateur. Elle devra notamment permettre :

\begin{itemize}
    \item La modification du contenu des cellules.
    \item La modification des lignes et des colonnes.
    \item La mise en forme des cellules.
    \item La modification des polices.
    \item La modification des couleurs.
    \item La modification des bordures.
    \item La modification des alignements.
    \item La gestion des formats de données pris en charge.
    \item La consultation et la modification des différentes feuilles.
\end{itemize}

Les fonctionnalités d'édition seront limitées aux besoins définis pour le projet et n'auront pas pour objectif de reproduire l'ensemble des fonctionnalités de Microsoft Excel ou Google Sheets.

\section{Conversion Excel vers Word}

L'application devra permettre de :

\begin{itemize}
    \item Sélectionner une feuille du classeur.
    \item Lancer la conversion de la feuille sélectionnée.
    \item Générer un document Word.
    \item Conserver la structure du tableau source.
    \item Conserver les bordures et les couleurs.
    \item Conserver les polices et les alignements pris en charge.
    \item Conserver les formats de données pris en charge.
    \item Adapter automatiquement la mise en page du document Word.
    \item Générer un document Word directement exploitable.
    \item Télécharger le document Word généré.
\end{itemize}

La conversion devra accorder une importance particulière à la fidélité de la présentation. Le document Word généré devra conserver au maximum la structure et la mise en forme de la feuille Excel afin d'éviter les déformations, débordements ou réajustements manuels.
